% !TeX root = ../../Book2.tex
% 纲要标题：## 附录 C　Azumaya 代数与非交换几何简介
% 纲要位置：工作记录/计划纲要.md，第 1824 行。

\chapter{Azumaya 代数与非交换几何简介}
\label{b2:app:C}

\begin{chapterabstract}
中心单代数在交换基环上的推广，需要用忠实有限生成投射模代替非零有限维向量空间。在底模已具备这些条件时，左右乘作用是否给出全部自同态，可以通过剩余域纤维检测；由此得到的 Azumaya 性在任意底环变换下保持。本附录证明这些局部判据，构造交换环的 Brauer 群并刻画其零类，再以向量丛的局部坐标与 Morita 等价解释自同态代数。实四元数与局部秩不同的投射模，分别说明在原底环上分裂及全局矩阵表示的限制。
\end{chapterabstract}

\begin{learninggoals}
\begin{enumerate}
\item 从左右乘作用识别 Azumaya 代数，区分它与单环、矩阵代数及分裂代数。
\item 运用有限投射模的局部自由性、剩余域判据和底环变换公式。
\item 证明 Brauer 等价关系与群运算的良定义，并解释零类的含义。
\item 将局部基、自同态代数的共轭变换以及模范畴等价对应到具体例子。
\end{enumerate}
\end{learninggoals}

在域上，\(M_n(k)\) 的矩阵单位提供全局坐标；换成交换环 \(R\)，一个有限投射模 \(P\) 可能只在各个局部环 \(R_{\mathfrak p}\) 上变成自由模。此时 \(\operatorname{End}_R(P)\) 在每个点附近都像矩阵代数，却未必能在整个环上用同一组矩阵单位书写。局部上相同的计算怎样拼回一个全局代数，成为域上理论移向基环时的新问题。

\ProseParagraphBreak

另一个变化出现在矩阵代数以外的中心单代数中：域上的除代数可能需要扩域才分裂，而换成交换环以后，还须同时处理模的局部自由性与代数的分裂。这里先要求代数的底模忠实、有限生成且投射，再检验各剩余域纤维是否中心单。在这个前提下，纤维条件恰好保证左右乘作用给出全部线性自同态，也就刻画了 Azumaya 代数。底模的有限投射性需要另外验证。

\ProseParagraphBreak
域上理论与交换基环上的对应如\cref{b2:tab:C-field-ring-comparison}所示。表中自同态代数是分裂代数的模型；一般中心单代数或 Azumaya 代数未必在原底域或底环上分裂。

\begin{table}[htbp]
\centering
\caption{域上理论与交换基环上的对应}
\label{b2:tab:C-field-ring-comparison}
\begin{tabularx}{\textwidth}{@{}XX@{}}
\toprule
域 \(k\) 上 & 非零交换环 \(R\) 上 \\
\midrule
非零有限维向量空间 & 忠实有限生成投射模 \\
分裂代数 \(M_n(k)=\operatorname{End}_k(k^n)\) & 分裂代数 \(\operatorname{End}_R(P)\) \\
有限维中心单代数 & Azumaya 代数 \\
全局基及换基矩阵 & 主开集上的局部基及粘合矩阵 \\
以矩阵代数为因子的稳定等价 & 以忠实有限投射模的自同态代数为因子的稳定等价 \\
\bottomrule
\end{tabularx}
\end{table}

有限生成与投射性使模在每个点附近具有有限自由基，忠实性再保证局部秩处处为正，基环没有被整个模遗漏的部分。这三项条件合在一起，才对应处处具有正秩的有限向量丛。本附录使用素理想谱、主开集、局部化及剩余域纤维的仿射代数描述；向量丛由给定投射模的局部坐标解释，不预设层与概形的一般理论。

\ProseParagraphBreak
本附录使用\cref{b2:prop:finite-projective-summand}与\cref{b2:prop:finite-dual-basis}中的有限投射模刻画、\cref{b2:thm:noncommutative-nakayama}、\cref{b2:thm:morita-equivalence}及\cref{b2:thm:csa-enveloping-isomorphism}。域上 Brauer 群的定义与除代数代表见\secref{b2:sec:1801}。

\ProseParagraphBreak
\secref{b2:sec:C01}从定义和基本例子建立对象，\secref{b2:sec:C02}给出局部检验方法。\secref{b2:sec:C03}构造 Brauer 群，并证明零类恰由分裂代数代表；\secref{b2:sec:C04}再说明局部基的粘合怎样作用于自同态代数，以及分裂代数为何与底环有等价的模范畴。

% !TeX root = ../../Book2.tex
% 纲要标题：### C.1 中心单代数的交换基环推广
% 纲要位置：工作记录/计划纲要.md，第 1826 行。

% 纲要要点（供目录核对）：
% * 用忠实有限生成投射模与左右乘包络映射推广中心单代数；分清单环、Azumaya、原基环分裂及矩阵代数
% #### C.1.1 Azumaya 代数的定义
% #### C.1.2 域情形与矩阵代数

\section{中心单代数的交换基环推广}
\label{b2:sec:C01}

中心单代数的定义包含域上的有限维性，但若把底域换成交换环，有限生成模未必有基，代数本身也通常不再是单环。例如 \(M_n(\mathbb Z)\) 有非零真双边理想 \(2M_n(\mathbb Z)\)，但在每个素数 \(p\) 处取剩余域纤维，仍得到中心单代数 \(M_n(\mathbb F_p)\)。所以推广的对象不能要求自身是单环，而要让各点上的中心单结构来自同一个有限投射模，并由左右乘作用统一描述。本附录的基环 \(R\) 交换、含幺且非零，代数的结构同态取值于中心；底环变换的目标环也约定非零。

\subsection{Azumaya 代数的定义}
\label{b2:subsec:C01:01}

对一个 \(R\) 代数 \(A\)，其反代数 \(A^{\mathrm{op}}\) 使用相反次序的乘法。左右乘给出 \(R\) 代数同态
\begin{equation}\label{b2:eq:azumaya-enveloping}
\mu_A:A\otimes_R A^{\mathrm{op}}\longrightarrow\operatorname{End}_R(A),
\qquad a\otimes b^{\mathrm{op}}\longmapsto(x\mapsto axb).
\end{equation}
例如，两次作用的复合把 \(x\) 送到 \(aa'xb'b\)，这正对应张量代数中的乘积 \(aa'\otimes(b'b)^{\mathrm{op}}\)。反代数保证右乘一侧的次序正确。

\begin{definition}\label{b2:def:azumaya-algebra}
设 \(R\) 为非零交换含幺环、\(A\) 为结合含幺 \(R\) 代数，其结构同态的像位于 \(A\) 的中心。若 \(A\) 满足以下两项条件，则称它为 \(R\) 上的\term[Azumaya-daishu]{Azumaya 代数}[Azumaya algebra]：
\begin{enumerate}
\item 作为 \(R\) 模，\(A\) 忠实、有限生成且投射。这里忠实指 \(rA=0\) 蕴含 \(r=0\)；有限生成且投射等价于存在一个 \(R\) 模 \(Q\) 及正整数 \(N\)，使 \(A\oplus Q\cong R^N\)。
\item 左右乘给出的 \(R\) 代数同态 \(\mu_A:A\otimes_RA^{\mathrm{op}}\to\operatorname{End}_R(A)\)、\(a\otimes b^{\mathrm{op}}\mapsto(x\mapsto axb)\) 是同构。
\end{enumerate}
\end{definition}
\symidx{\(\mu_A\)：Azumaya 代数的左右乘作用映射}

域上\cref{b2:thm:csa-enveloping-isomorphism}正是用同一个包络映射刻画中心单代数；这里将底域换成交换环，并用忠实有限生成投射模代替非零有限维向量空间。定义中的直和刻画由\cref{b2:prop:finite-projective-summand}保证。它允许模没有全局基，却仍能在局部选择有限基。忠实性则使结构同态 \(R\rightarrow A\) 单射：若 \(r1_A=0\)，就有 \(rA=0\)。因此可以把 \(R\) 看作 \(A\) 中的标量。\cref{b2:prop:azumaya-center}将进一步证明，\(A\) 的中心恰好是这些标量。

\ProseParagraphBreak
\(R\) 本身是 \(R\) 上的 Azumaya 代数，因为 \(R\otimes_R R\cong R\cong\operatorname{End}_R(R)\)，且这些同构把 \(\mu_R\) 变成恒等映射。若 \(0\ne I\subsetneq R\)，商环 \(R/I\) 的标量作用却零化了整个 \(I\)，不能作为同一个基环 \(R\) 上的 Azumaya 代数。例如取 \(R=k\times k\)、\(I=0\times k\)，则 \(A=R/I\cong k\) 是有限投射 \(R\) 模，包络映射也同构，却完全遗漏了基环的第二个分量。这说明忠实性是独立需要的条件。改变底环后，\(R/I\) 作为自身上的代数满足定义；底环是概念的一部分。

\subsection{域情形与矩阵代数}
\label{b2:subsec:C01:02}

\begin{proposition}\label{b2:prop:azumaya-field-matrix}
域 \(k\) 上的 Azumaya 代数恰为有限维中心单 \(k\) 代数。对任意非零交换环 \(R\) 及 \(n\ge1\)，\(M_n(R)\) 都是 \(R\) 上的 Azumaya 代数。
\end{proposition}
\begin{proof}
域上的忠实有限生成投射模恰为非零有限维向量空间。在此前提下，\cref{b2:thm:csa-enveloping-isomorphism}说明，左右乘作用同构等价于代数中心单。这给出第一个断言的两个方向。

对 \(A=M_n(R)\)，矩阵单位 \(E_{ij}\) 给出 \(A\) 的 \(R\) 基，所以它忠实且有限自由。直接计算得
\[
\mu_A(E_{ij}\otimes E_{\ell m}^{\mathrm{op}})(E_{pq})
=\delta_{jp}\delta_{q\ell}E_{im}.
\]
随着 \(i,j,\ell,m\) 变化，右侧恰为自由模 \(A\cong R^{n^2}\) 的自同态环中的标准矩阵单位：它把一个指定基向量 \(E_{j\ell}\) 送到一个指定基向量 \(E_{im}\)，并把其余基向量送到零。因此 \(\mu_A\) 把来源的一组基双射到目标的一组基，故为同构。
\end{proof}

例如 \(M_2(\mathbb Z/4\mathbb Z)\) 同样满足定义：基环可以有零因子，矩阵大小也不必在基环中可逆。Azumaya 性不是把矩阵迹除以矩阵大小所得到的性质。

\ProseParagraphBreak
另一端的例子是实四元数代数 \(\mathbb H\)。\cref{b2:prop:quaternion-division-or-split}及范数公式\eqref{b2:eq:quaternion-norm}说明，它是中心除 \(\mathbb R\) 代数：非零四元数的范数是四个实数平方之和，因而不为零。所以 \(\mathbb H\) 是 \(\mathbb R\) 上的 Azumaya 代数，却不同构于 \(M_2(\mathbb R)\)，因为后者有非零零因子。定义允许这种在原基域上未分裂的中心单代数；要求 Azumaya 性，并没有要求原基环上已经存在矩阵单位。

% !TeX root = ../../Book2.tex
% 纲要标题：### C.2 Azumaya 代数的局部性质
% 纲要位置：工作记录/计划纲要.md，第 1834 行。

% 纲要要点（供目录核对）：
% * 局部化与剩余域纤维分别检验；Nakayama、目标投射性和迹单位的使用位置，主开有限覆盖、任意变基、中心、泛纤维反例与单位判据；平展正向提升明确外部归属，反向完整下降
% #### C.2.1 有限投射模的局部检验
% #### C.2.2 纤维判据与底环变换

\section{Azumaya 代数的局部性质}
\label{b2:sec:C02}

当代数的底模已知忠实、有限生成且投射时，左右乘作用是两个有限投射模之间的映射，因而可以在剩余域上检验。本节建立这一检验所需的局部模论。这里的局部化使用\cref{b2:def:localized-module}的分式构造，零判据为式\eqref{b2:eq:localization-zero}；局部化保持正合性，来自\cref{b2:prop:localization-flat}与其张量描述。

\subsection{有限投射模的局部检验}
\label{b2:subsec:C02:01}

设 \(R\) 为非零交换含幺环。\(R\) 的全部素理想组成的集合称为\term[sulixiang-pu]{素理想谱}[prime spectrum]，记为 \(\operatorname{Spec}R\)。对 \(\mathfrak p\in\operatorname{Spec}R\)，令 \(R_{\mathfrak p}=(R\setminus\mathfrak p)^{-1}R\)，并对 \(R\) 模 \(M\) 令 \(M_{\mathfrak p}=R_{\mathfrak p}\otimes_R M\)。环 \(R_{\mathfrak p}\) 的唯一极大理想为 \(\mathfrak pR_{\mathfrak p}\)：分子不在 \(\mathfrak p\) 中的分式可逆，分子在 \(\mathfrak p\) 中的分式都属于这个真理想。它的\term[shengyu-yu]{剩余域}[residue field]及 \(M\) 在该点的\term[xianwei-mo]{纤维}[fiber of a module]分别记为
\[
\kappa(\mathfrak p)=R_{\mathfrak p}/\mathfrak pR_{\mathfrak p},
\qquad M(\mathfrak p)=\kappa(\mathfrak p)\otimes_R M.
\]
若 \(\mathfrak m\) 极大，则 \(\kappa(\mathfrak m)=R/\mathfrak m\)。这里把素理想看作一个点，剩余域是在该点取值所用的系数域，而纤维是这个域上的向量空间。局部化保留了局部环上的模，取纤维还要模去其极大理想：
\[
M(\mathfrak p)\cong M_{\mathfrak p}/\mathfrak pM_{\mathfrak p}.
\]
\symidx{\(\operatorname{Spec}R\)：交换环的素理想谱}
\symidx{\(\kappa(\mathfrak p)\)：素理想处的剩余域}
\symidx{\(M(\mathfrak p)\)：模在素理想处的纤维}

\begin{lemma}[局部消失与有限投射模]\label{b2:lem:azumaya-local-toolkit}
设 \(R\) 为非零交换含幺环。对 \(R\) 模 \(X\) 和素理想 \(\mathfrak p\)，记 \(R_{\mathfrak p}=(R\setminus\mathfrak p)^{-1}R\)、\(X_{\mathfrak p}=R_{\mathfrak p}\otimes_RX\)、\(\kappa(\mathfrak p)=R_{\mathfrak p}/\mathfrak pR_{\mathfrak p}\)，并记纤维 \(X(\mathfrak p)=\kappa(\mathfrak p)\otimes_RX\)。
\begin{enumerate}
\item 对任意 \(R\) 模 \(M\)，\(M=0\) 当且仅当每个极大理想 \(\mathfrak m\) 处的 \(M_{\mathfrak m}=0\)。若 \(M\) 有限生成，还可等价地要求全部纤维 \(M(\mathfrak m)=0\)。
\item 对任意有限生成投射 \(R\) 模 \(P\)，每个 \(P_{\mathfrak p}\) 都是有限自由 \(R_{\mathfrak p}\) 模。
\item 若 \(P\) 为有限生成 \(R\) 模、\(Q\) 为有限生成投射 \(R\) 模，且 \(u:P\rightarrow Q\) 为 \(R\) 模同态，则 \(u\) 为同构，当且仅当它在每个极大理想 \(\mathfrak m\) 处诱导的 \(u(\mathfrak m):P(\mathfrak m)\rightarrow Q(\mathfrak m)\) 为同构。
\end{enumerate}
\end{lemma}
\begin{proof}
若 \(0\ne x\in M\)，要找一个局部化仍能看见它的点，只需让分母避开它的零化子。具体地，\(I=\{r\in R:rx=0\}\) 是真理想，取包含它的极大理想 \(\mathfrak m\)。若 \(x/1=0\) 于 \(M_{\mathfrak m}\)，分式零判据给出 \(s\notin\mathfrak m\) 且 \(sx=0\)，与 \(I\subseteq\mathfrak m\) 矛盾。故全部局部化为零时 \(M=0\)，这一步不需要有限生成。若进一步知道 \(M\) 有限生成且 \(M(\mathfrak m)=0\)，就有 \(M_{\mathfrak m}=\mathfrak mR_{\mathfrak m}M_{\mathfrak m}\)。对局部环 \(R_{\mathfrak m}\)，Jacobson 根是其唯一极大理想；\cref{b2:thm:noncommutative-nakayama}因而使有限生成模 \(M_{\mathfrak m}\) 为零。这里从纤维消失回到局部化消失，才用到有限生成。反向都由取张量立即得到。

设 \(T\) 是局部环，极大理想为 \(\mathfrak n\)，而 \(L\) 为有限生成投射 \(T\) 模。我们想把剩余域上的基提升成 \(L\) 的基：提升 \(L/\mathfrak nL\) 的一组基，得到 \(f:T^r\rightarrow L\)。\cref{b2:cor:nakayama-generators}先保证这些提升生成 \(L\)，所以 \(f\) 满射。要进一步排除它们之间的关系，利用投射性写 \(T^r\cong L\oplus K\)，其中 \(K\) 作为有限自由模的直和因子仍有限生成。模去 \(\mathfrak n\) 后，\(f\) 已是同构，故 \(K/\mathfrak nK=0\)；再由 Nakayama 引理，\(K=0\)。于是 \(L\cong T^r\)。把 \(P\oplus P'\cong R^N\) 局部化可知 \(P_{\mathfrak p}\) 仍有限生成且投射，取 \(T=R_{\mathfrak p}\) 即得第二项。

对第三项，先检测满射性。\(\operatorname{coker}u\) 是有限生成模 \(Q\) 的商模，张量的右正合性使它的全部极大纤维为零。第一项说明 \(u\) 满射。要检测单射性，不能直接假定取纤维保持核；这里需要目标 \(Q\) 的投射性，使该满射分裂，得到 \(P\cong\ker u\oplus Q\)。核作为有限生成模 \(P\) 的直和因子仍有限生成；对分裂列取纤维仍分裂，纤维同构假设使核的全部纤维为零。再次应用第一项，得到核为零。反向由任何函子保持同构得到。
\end{proof}

\begin{lemma}[有限投射模与底环变换]\label{b2:lem:azumaya-projective-base-change}
设 \(R,S\) 为非零交换含幺环、\(R\to S\) 为环同态，\(P,Q\) 为有限生成投射 \(R\) 模，则 \(P,Q\) 的对偶模和张量积仍有限生成且投射，并有自然同构
\begin{align}
S\otimes_R\operatorname{Hom}_R(P,Q)
&\cong\operatorname{Hom}_S(S\otimes_RP,S\otimes_RQ),\label{b2:eq:azumaya-hom-base-change}\\
\operatorname{End}_R(P)\otimes_R\operatorname{End}_R(Q)
&\cong\operatorname{End}_R(P\otimes_RQ).\label{b2:eq:azumaya-end-tensor}
\end{align}
后一个映射把 \(f\otimes g\) 送到 \(p\otimes q\mapsto f(p)\otimes g(q)\)，并为代数同构。若 \(P\) 还忠实，则它是生成元，且 \(S\otimes_RP\) 是忠实有限生成投射 \(S\) 模。两个忠实有限生成投射模的张量积也忠实。
\end{lemma}
\begin{proof}
有限自由模直和分解取对偶、取底环变换或逐因子取张量后仍然分裂，所以所列模均为有限自由模的直和因子。\cref{b2:prop:finite-projective-hom-tensor}给出 \(\operatorname{Hom}_R(P,Q)\cong P^*\otimes_RQ\)。对偶的底环变换公式先在有限自由模上由坐标成立，再限制到直和因子，得到式\eqref{b2:eq:azumaya-hom-base-change}。式\eqref{b2:eq:azumaya-end-tensor}在极大理想处化为两个有限自由模的矩阵单位对应，故由\cref{b2:lem:azumaya-local-toolkit}的同构判据成立；映射的公式直接保留乘法。

现在设 \(P\) 忠实。若某个 \(P_{\mathfrak m}=0\)，取 \(P\) 的有限生成组，分别用不在 \(\mathfrak m\) 中的元素零化它们，再取这些元素的乘积，得到 \(s\notin\mathfrak m\) 且 \(sP=0\)。忠实性使 \(s=0\)，矛盾。因此每个 \(P_{\mathfrak m}\) 都自由且秩为正。

回忆\cref{b2:prop:trace-generator-criterion}中的迹理想
\[
I=\operatorname{tr}_R(P)=\sum_{\lambda\in P^*}\lambda(P)\subseteq R,
\qquad P^*=\operatorname{Hom}_R(P,R).
\]
它由全部 \(R\) 线性泛函在 \(P\) 上的取值生成。式\eqref{b2:eq:azumaya-hom-base-change}使取局部对偶与取局部化相容，故 \(I_{\mathfrak m}\) 由 \(P_{\mathfrak m}\) 的对偶泛函在模元素上的取值生成。由于局部自由秩为正，选一个基向量及其坐标泛函，便有 \(I_{\mathfrak m}=R_{\mathfrak m}\)。因此 \((R/I)_{\mathfrak m}=0\) 对全部 \(\mathfrak m\) 成立，局部消失判据给出 \(I=R\)。虽然定义迹理想时允许所有泛函，表示其中的单位元只用有限多个取值；将系数吸收到泛函中，可写
\begin{equation}\label{b2:eq:azumaya-trace-unit}
1=\sum_{i=1}^t\lambda_i(p_i).
\end{equation}
由\cref{b2:prop:trace-generator-criterion}，这就说明 \(P\) 为生成元。

底环变换后式\eqref{b2:eq:azumaya-trace-unit}仍成立。若 \(s\in S\) 零化 \(S\otimes_RP\)，对这个等式乘以 \(s\)，再利用各扩张泛函的 \(S\) 线性性，得到 \(s=0\)。因此底环变换后仍忠实。若 \(Q\) 也忠实，将两者表示单位元的等式相乘，便得到 \(P\otimes_RQ\) 上相应的等式；同一论证证明张量积忠实。
\end{proof}

忠实性在这里与有限投射性一起使用：忠实性与有限生成共同保证每个局部化非零，投射性再使局部对偶中有坐标泛函，从而得到式\eqref{b2:eq:azumaya-trace-unit}。单有忠实性既不能推出模为生成元，也不能保证任意底环变换后仍忠实；\(\mathbb Q\) 作为 \(\mathbb Z\) 模的例子见\subsecref{b2:subsec:1601:01}及下文\cref{b2:prop:C-fiber-hypothesis-boundaries}。

\ProseParagraphBreak
局部环上的自由基还可以延伸到点的一个邻域。为说明这里的邻域，给前面的素理想集合加上拓扑：对 \(f\in R\)，记 \(D(f)=\{\mathfrak p\in\operatorname{Spec}R\mid f\notin\mathfrak p\}\)。这些集合满足 \(D(1)=\operatorname{Spec}R\)、\(D(f)\cap D(g)=D(fg)\)；以它们为开集的一组基，得到素理想谱上的\term[Zariski-tuopu]{Zariski 拓扑}[Zariski topology]。集合 \(D(f)\) 称为主开集，在其中允许 \(f\) 可逆，对应于底环局部化 \(R_f=R[1/f]\)。
\symidx{\(D(f)\)：素理想谱的主开集}

\begin{proposition}[局部基的延伸]\label{b2:prop:projective-basic-open-basis}
设 \(R\) 为非零交换含幺环、\(P\) 为有限生成投射 \(R\) 模。记 \(\operatorname{Spec}R\) 为 \(R\) 的素理想集合；对 \(f\in R\)，记 \(D(f)=\{\mathfrak p\in\operatorname{Spec}R\mid f\notin\mathfrak p\}\)、\(R_f=R[1/f]\) 和 \(P_f=R_f\otimes_RP\)。每个 \(\mathfrak p\in\operatorname{Spec}R\) 都有一个主开邻域 \(D(f)\)，使 \(P_f\) 为有限自由 \(R_f\) 模。这样的主开集可以选有限多个覆盖 \(\operatorname{Spec}R\)。
\end{proposition}
\begin{proof}
\cref{b2:lem:azumaya-local-toolkit}说明 \(P_{\mathfrak p}\) 自由。其有限个基向量是 \(P\) 中元素的分式，把每个基向量乘以一个可逆分母后，仍得到一组基，故可选 \(p_1,\ldots,p_r\in P\)，使映射 \(\alpha:R^r\rightarrow P\) 在 \(\mathfrak p\) 处为同构。

由式\eqref{b2:eq:azumaya-hom-base-change}，局部逆映射可写成某个 \(\beta\in\operatorname{Hom}_R(P,R^r)\) 除以 \(s\notin\mathfrak p\)。要使它在一个主开集上已经成为逆映射，只需同时消去两个复合与恒等映射之间的误差。在 \(R_s\) 上，考察 \((\beta/s)\alpha-1\) 与 \(\alpha(\beta/s)-1\)，它们在 \(\mathfrak p\) 处为零。分别取 \(R_s^r\) 和 \(P_s\) 的有限生成组，对这有限多个值应用分式零判据，再把所得分母相乘，可以选 \(t\notin\mathfrak p\)，使两映射在 \(R_{st}\) 上都为零。因此 \(f=st\) 满足要求。有限基、局部逆映射和有限多个误差值，正是这里能够一次清去分母的有限数据。

若所有这些 \(f\) 生成的理想为真理想，取包含它的极大理想，就得到一个不在任何 \(D(f)\) 中的素理想，矛盾。因此有有限个 \(f_1,\ldots,f_u\) 及 \(a_i\in R\)，使 \(\sum_i a_if_i=1\)。任何素理想不能同时包含这些 \(f_i\)，故对应的有限个主开集已经覆盖全谱。
\end{proof}

最后这个单位理想论证对任意主开覆盖都成立。因为主开集构成拓扑的一组基，它也说明 \(\operatorname{Spec}R\) 是拟紧空间。

\subsection{纤维判据与底环变换}
\label{b2:subsec:C02:02}

设 \(R\) 为非零交换含幺环、\(A\) 为结合含幺 \(R\) 代数，且结构同态取值于中心。对 \(\mathfrak p\in\operatorname{Spec}R\)，局部化与取纤维分两步进行：
\[
\begin{tikzcd}[column sep=3.5em,row sep=2em]
R\arrow[r]\arrow[d]&R_{\mathfrak p}\arrow[r,two heads]\arrow[d]&\kappa(\mathfrak p)\arrow[d]\\
A\arrow[r]&A_{\mathfrak p}\arrow[r,two heads]&A(\mathfrak p)
\end{tikzcd}
\]
这里 \(A_{\mathfrak p}=R_{\mathfrak p}\otimes_RA\) 为局部化代数，\(A(\mathfrak p)=\kappa(\mathfrak p)\otimes_RA\) 为纤维代数；竖箭头为结构同态，两个方块均交换。第一步使 \(R\setminus\mathfrak p\) 中的元素可逆，第二步再模去 \(\mathfrak pR_{\mathfrak p}\)，因而
\[
A(\mathfrak p)\cong A_{\mathfrak p}/\mathfrak pA_{\mathfrak p}.
\]
局部化代数的底环仍是局部环 \(R_{\mathfrak p}\)，纤维代数的底环则是剩余域 \(\kappa(\mathfrak p)\)。纤维由代数自身作为模的商构造得到；即使代数嵌在矩阵代数中，也不能未经检验就把它等同于逐项代入剩余域所得的矩阵像，具体区别见\cref{b2:prop:C-generic-fiber-insufficient}。

\begin{theorem}[Azumaya 代数的纤维判据]\label{b2:thm:azumaya-fiber-criterion}
设 \(R\) 为非零交换含幺环、\(A\) 为结合含幺 \(R\) 代数，其结构同态取值于中心，且 \(A\) 作为 \(R\) 模忠实、有限生成且投射。对素理想 \(\mathfrak p\)，记 \(R_{\mathfrak p}=(R\setminus\mathfrak p)^{-1}R\)、\(A_{\mathfrak p}=R_{\mathfrak p}\otimes_RA\)、\(\kappa(\mathfrak p)=R_{\mathfrak p}/\mathfrak pR_{\mathfrak p}\)、\(A(\mathfrak p)=\kappa(\mathfrak p)\otimes_RA\)。以下条件等价：
\begin{enumerate}
\item \(A\) 为 Azumaya \(R\) 代数；
\item 对每个极大理想 \(\mathfrak m\)，\(A(\mathfrak m)\) 为中心单 \(\kappa(\mathfrak m)\) 代数；
\item 对每个素理想 \(\mathfrak p\)，\(A_{\mathfrak p}\) 为 Azumaya \(R_{\mathfrak p}\) 代数。
\end{enumerate}
此外，若 \(A\) 满足这些条件，且 \(R\to S\) 为非零交换含幺环的同态，则 \(S\otimes_RA\) 为 Azumaya \(S\) 代数。
\end{theorem}
\begin{proof}
Azumaya 条件要检验的是 \(\mu_A\) 是否同构，而纤维上的中心单条件正好检验它的底环变换。\cref{b2:lem:azumaya-projective-base-change}保证 \(A\otimes_RA^{\mathrm{op}}\) 与 \(\operatorname{End}_R(A)\) 都有限生成且投射，并使 \(\mu_A\) 在 \(\kappa(\mathfrak m)\) 上的底环变换恰为 \(\mu_{A(\mathfrak m)}\)。该引理的忠实性论证还给出 \(A_{\mathfrak m}\) 的自由秩为正，所以 \(A(\mathfrak m)\ne0\)；仅有投射性并不能保证这一点。由\cref{b2:prop:azumaya-field-matrix}，该纤维映射同构等价于纤维代数中心单；再由\cref{b2:lem:azumaya-local-toolkit}检测同构，第一、二项等价。

若第一项成立，对任意 \(R\rightarrow S\) 取底环变换。\cref{b2:lem:azumaya-projective-base-change}给出忠实有限投射性，又有
\[
(S\otimes_RA)\otimes_S(S\otimes_RA)^{\mathrm{op}}
\cong S\otimes_R(A\otimes_RA^{\mathrm{op}}),
\]
以及自同态代数的底环变换同构。在这些识别下，\(\mu_{S\otimes_RA}\) 是 \(S\otimes_R\mu_A\)，故为同构。这里不要求 \(S\) 平坦：任何函子都保持同构，而有限投射性保证所用的自同态变基公式，迹理想中的单位元等式保证变基后的忠实性。这证明底环变换断言，也证明第一项蕴含第三项。若第三项成立，再对 \(R_{\mathfrak m}\rightarrow\kappa(\mathfrak m)\) 作底环变换便得第二项，等价性证毕。
\end{proof}

判据要求检验所有极大纤维；只在分式域上或某个主开集上成为矩阵代数，仍不足以推出原代数 Azumaya，见\cref{b2:prop:C-generic-fiber-insufficient}。局部模的同构判据还可检测有限投射代数中单个元素的可逆性，见\cref{b2:prop:C-unit-fiber-criterion}。

\begin{corollary}\label{b2:cor:azumaya-endomorphism}
设 \(R\) 为非零交换含幺环。若 \(P\) 是忠实有限生成投射 \(R\) 模，则 \(\operatorname{End}_R(P)\) 是 Azumaya \(R\) 代数。
\end{corollary}
\begin{proof}
由\cref{b2:lem:azumaya-projective-base-change}，自同态模为有限投射模。若 \(r\) 零化所有自同态，它就零化恒等映射，故零化 \(P\)，于是 \(r=0\)。每个极大纤维由式\eqref{b2:eq:azumaya-hom-base-change}同构于非零有限维向量空间的自同态代数，所以是中心单代数。纤维判据给出结论。
\end{proof}

\begin{proposition}\label{b2:prop:azumaya-center}
设 \(R\) 为非零交换含幺环、\(A\) 为 Azumaya \(R\) 代数，则 \(A\) 的中心为 \(R1_A\)。此外，对任意忠实有限生成投射 \(R\) 模 \(P\)，有
\[
Z\bigl(\operatorname{End}_R(P)\bigr)
=\{r\operatorname{id}_P\mid r\in R\}.
\]
这两处的标量表示都唯一。
\end{proposition}
\begin{proof}
先设 \(P\) 为忠实有限生成投射 \(R\) 模，\(T\in\operatorname{End}_R(P)\) 与全部自同态交换。取式\eqref{b2:eq:azumaya-trace-unit}中的 \(p_i,\lambda_i\)。要证明 \(T(x)\) 是同一个标量乘以 \(x\)，可先乘上这些取值后再求和；为把 \(T\) 从 \(x\) 移到 \(p_i\)，让它与自同态 \(y\mapsto\lambda_i(y)x\) 交换，再在 \(p_i\) 处取值，得到
\[
\lambda_i(p_i)T(x)=\lambda_i\bigl(T(p_i)\bigr)x.
\]
求和说明 \(T(x)=rx\)，其中 \(r=\sum_i\lambda_i\bigl(T(p_i)\bigr)\in R\) 不依赖于 \(x\)。反过来，标量乘法与每个 \(R\) 线性自同态交换；忠实性保证不同标量给出不同自同态。由此得到关于自同态代数中心的断言。

若 \(z\in Z(A)\)，左乘 \(z\) 与所有左右乘算子交换。\(\mu_A\) 满射，因而它与全部 \(R\) 线性自同态交换。取 \(P=A\)，由刚证结论有 \(L_z=r\operatorname{id}_A\)，在 \(1_A\) 处取值得 \(z=r1_A\)。反向包含是 \(R\) 代数的约定，而忠实性保证这种标量表示唯一。
\end{proof}

纤维判据中的“中心单”不能换成“在剩余域上已经分裂”而仍声称二者等价。对基环 \(\mathbb R\)，唯一素理想为零，局部化和剩余域都不改变基环；\(\mathbb H\) 的纤维仍为除代数。换到 \(\mathbb C\) 后它才成为 \(M_2(\mathbb C)\)。因此 Azumaya 性并不意味着在 Zariski 局部化后就能选出矩阵单位。

\ProseParagraphBreak
这个例子也提示了分裂所需的另一种局部操作：除了使指定元素可逆，还允许剩余域上的有限可分扩张。平展代数同时包含这两种操作。本附录用以下仿射条件描述它：设 \(R\to S\) 为交换含幺环的同态，若 \(S\) 作为交换 \(R\) 代数有有限个生成元和有限条多项式关系，作为 \(R\) 模平坦，且对每个素理想 \(\mathfrak p\)，\(\kappa(\mathfrak p)\otimes_RS\) 是有限个有限可分 \(\kappa(\mathfrak p)\) 域扩张的直积，则称 \(S\) 为\term[pingzhan-daishu]{平展代数}[étale algebra]。域上的有限可分扩张与底环的主开局部化都满足这些条件。

\ProseParagraphBreak
这里允许零纤维，即空直积。例如 \(\kappa(\mathfrak p)\otimes_RR_f\cong\kappa(\mathfrak p)_f\)，若 \(f\in\mathfrak p\)，它要求把剩余域中的零变成单位，因而为零环。这表示 \(\operatorname{Spec}R_f\) 没有点映到 \(\mathfrak p\)，即这个主开集没有覆盖该点。因此下面除了要求各处平展分裂，还必须要求这些谱映射覆盖全谱。

\begin{theorem}[平展局部分裂]\label{b2:thm:azumaya-etale-splitting}
设 \(R\) 为非零交换含幺环、\(A\) 为结构同态取值于中心的结合含幺 \(R\) 代数，且 \(A\) 作为 \(R\) 模忠实、有限生成且投射。则 \(A\) 是 Azumaya \(R\) 代数，当且仅当存在有限族非零交换平展 \(R\) 代数 \(S_1,\ldots,S_t\)，满足以下两项条件：
\begin{enumerate}
\item 每个素理想 \(\mathfrak p\in\operatorname{Spec}R\) 都是某个 \(S_i\) 的素理想在 \(R\) 中的收缩，即这些谱映射的像联合覆盖 \(\operatorname{Spec}R\)。
\item 对每个 \(i\)，存在正整数 \(n_i\) 和 \(S_i\) 代数同构
\[
S_i\otimes_RA\cong M_{n_i}(S_i).
\]
\end{enumerate}
\end{theorem}
\begin{proofsketch}
正向需要把纤维上的可分分裂代数提升到点附近，这一提升不由前面的局部模论直接给出。这里只说明所调用的外部结论及其用途：每个点都有开邻域 \(U\) 以及有限平展满射 \(U'\to U\)，使 \(A\) 在 \(U'\) 上成为矩阵代数。这一提升及分裂构造见\cite[Théorème 5.1, Proposition 5.4, Théorème 5.7]{Grothendieck1966BrauerI}；此处省去提升步骤的证明。把 \(U\) 缩小为主开集，\(U'\) 仍为仿射，从而得到所需的平展 \(R\) 代数。\(\operatorname{Spec}R\) 的拟紧性使覆盖可以取为有限族。

反向可以用前面的模论完整证明。令 \(S=\prod_i S_i\)，各 \(S_i\) 平坦，有限直积 \(S\) 也平坦。谱覆盖还保证它的张量函子检测非零模。确实，对 \(0\ne x\in M\)，令 \(I=\operatorname{Ann}_R(x)\)，并选包含 \(I\) 的极大理想 \(\mathfrak m\)。覆盖条件给出某个 \(S_j\) 的素理想 \(\mathfrak q\)，收缩为 \(\mathfrak m\)；于是 \(IS_j\subseteq\mathfrak q\)，故 \(S_j\otimes_RRx\cong S_j/IS_j\ne0\)。平坦性使这个模嵌入 \(S_j\otimes_RM\)，所以 \(S\otimes_RM\ne0\)。因此 \(S\) 不仅平坦，而且张量后一个模为零会迫使原模为零，这就是这里所需的忠实平坦性。

在每个 \(S_i\) 上，\(\mu_A\) 经底环变换成为矩阵代数的左右乘同构。由式\eqref{b2:eq:azumaya-hom-base-change}，\(S\otimes_R\mu_A\) 因而为同构。平坦性使取张量保持核、余核，所以
\[
S\otimes_R\ker\mu_A=0,
\qquad S\otimes_R\operatorname{coker}\mu_A=0.
\]
忠实平坦性使 \(\ker\mu_A=0\) 且 \(\operatorname{coker}\mu_A=0\)。因此 \(\mu_A\) 同构，得到 Azumaya 性。
\end{proofsketch}

有限族指覆盖中代数的个数有限，不要求每个 \(S_i\) 都是有限 \(R\) 模。矩阵大小也可随基谱的不同部分变化。平展局部分裂是 Azumaya 代数的几何刻画；它允许像 \(\mathbb R\to\mathbb C\) 这样的可分扩张，因而比只取主开集的 Zariski 局部分裂更一般。

\begin{proposition}[纤维检验的假设边界]\label{b2:prop:C-fiber-hypothesis-boundaries}
非零 \(\mathbb Z\) 模 \(\mathbb Q\) 忠实，且对每个素数 \(p\)，有
\[
\mathbb Q_{(p)}\cong\mathbb Q\ne0,
\qquad
\mathbb F_p\otimes_{\mathbb Z}\mathbb Q=0.
\]
因而不能对任意模用全部极大纤维的消失检测原模的消失，也不能对任意忠实模断言底环变换后仍忠实。另设 \(k\) 为任意域、\(R=k[\varepsilon]/(\varepsilon^2)\)，则商映射 \(\rho:R\to k=R/(\varepsilon)\) 的唯一极大纤维为同构，但 \(\ker\rho=(\varepsilon)\ne0\)，且目标 \(k\) 不是投射 \(R\) 模。
\end{proposition}
\begin{proof}
非零整数不能零化 \(\mathbb Q\)，所以它作为 \(\mathbb Z\) 模忠实。任意不在 \((p)\) 中的整数在 \(\mathbb Q\) 上已经可逆，局部化的映射 \(q/s\mapsto q/s\) 给出 \(\mathbb Q_{(p)}\cong\mathbb Q\)。另一方面，\(p\mathbb Q=\mathbb Q\)，故
\[
\mathbb F_p\otimes_{\mathbb Z}\mathbb Q
\cong\mathbb Q/p\mathbb Q=0.
\]
这不违反\cref{b2:lem:azumaya-local-toolkit}，因为 \(\mathbb Q\) 并非有限生成 \(\mathbb Z\) 模：有限多个有理数的分母有公共倍数 \(d\)，它们的整数线性组合都在 \(d^{-1}\mathbb Z\) 中，不能包含 \(1/(2d)\)。所得零模在非零环 \(\mathbb F_p\) 上也不忠实。

对双数环 \(R\)，元素 \(a+b\varepsilon\) 在 \(a\ne0\) 时可逆，其逆为 \(a^{-1}-a^{-2}b\varepsilon\)，而 \((\varepsilon)\) 的商为域 \(k\)，所以 \((\varepsilon)\) 是唯一极大理想。对 \(\rho\) 在这个点取纤维，得到 \(k\otimes_RR\to k\otimes_Rk\)，即 \(k\) 上的恒等映射；原映射却有非零核 \((\varepsilon)\)。若 \(k\) 投射，\(\rho\) 就有 \(R\) 线性截面 \(s:k\to R\)。从 \(\rho(s(1))=1\) 可写 \(s(1)=1+b\varepsilon\)，但 \(k\) 被 \(\varepsilon\) 零化，线性性要求
\[
0=s(\varepsilon\cdot1)=\varepsilon s(1)=\varepsilon,
\]
矛盾。因此，纤维上的同构没有排除原核，正是目标投射性的缺失所允许的现象。
\end{proof}

\begin{proposition}[泛纤维不足以检验 Azumaya 性]\label{b2:prop:C-generic-fiber-insufficient}
设 \(k\) 为任意域、\(R=k[t]\)，并令
\[
A=\left\{
\begin{pmatrix}a&b\\tc&d\end{pmatrix}
\ \middle|\ a,b,c,d\in R
\right\}\subseteq M_2(R).
\]
代数 \(A\) 含幺，底模忠实且自由，秩为 \(4\)。它满足 \(A[1/t]\cong M_2(R[1/t])\)，从而泛纤维 \(k(t)\otimes_RA\cong M_2(k(t))\)；但在极大理想 \((t)\) 处的纤维 \(A/tA\) 不是单代数，所以 \(A\) 不是 Azumaya \(R\) 代数。将矩阵条目在 \(t=0\) 处求值的映射在该纤维上并不单射。
\end{proposition}
\begin{proof}
矩阵乘法保持左下条目为 \(t\) 的倍数，且 \(1_A=E_{11}+E_{22}\)，所以 \(A\) 确为含幺子代数。令
\[
e=E_{11},\qquad f=E_{22},\qquad x=E_{12},\qquad y=tE_{21}.
\]
每个元素唯一写成 \(ae+df+bx+cy\)，故这四个元素为 \(R\) 基，底模自由且忠实。使 \(t\) 可逆后，\(t^{-1}y=E_{21}\)，四个矩阵单位均在局部化中，因而 \(A[1/t]=M_2(R[1/t])\)。再把底环扩张到 \(k(t)\) 即得泛纤维的断言。

由于 \(A\) 的这四个元素构成自由基，它们的商类 \(\bar e,\bar f,\bar x,\bar y\) 构成 \(A/tA\) 的 \(k\) 基，特别是 \(\bar y\ne0\)。在 \(A\) 中有
\[
x^2=y^2=0,\qquad xy=te,\qquad yx=tf.
\]
又有 \(ex=x=xf\)、\(fy=y=ye\)，其余用 \(e,f\) 从两侧乘 \(x,y\) 的结果为零。因此在 \(A/tA\) 中，\(J=k\bar x+k\bar y\) 是非零真双边理想，且 \(J^2=0\)。纤维不是单代数，\cref{b2:thm:azumaya-fiber-criterion}便排除了 \(A\) 的 Azumaya 性。

把 \(A\) 中的矩阵条目在 \(t=0\) 处求值，得到满射到上三角矩阵代数的映射。它零化 \(tA\)，所以诱导 \(A/tA\) 上的映射，却把 \(\bar y\) 送到零。内在纤维仍保留基向量 \(y\) 的非零商类；外在的条目求值丢掉了它，故这两个构造不能混同。
\end{proof}

\begin{proposition}[有限投射代数中的单位检验]\label{b2:prop:C-unit-fiber-criterion}
设 \(R\) 为非零交换含幺环、\(A\) 为结构同态取值于中心的结合含幺 \(R\) 代数，且 \(A\) 作为 \(R\) 模有限生成且投射。对 \(a\in A\)，\(a\) 是 \(A\) 的单位，当且仅当对每个极大理想 \(\mathfrak m\)，它在纤维代数 \(A(\mathfrak m)=(R/\mathfrak m)\otimes_RA\) 中的像都是单位。这里不要求 \(A\) 忠实或满足 Azumaya 条件。
\end{proposition}
\begin{proof}
若 \(ab=ba=1_A\)，取任一纤维后仍有这两个等式，故正向成立。反之，左乘映射 \(L_a:A\to A\) 是 \(R\) 线性映射；在各极大纤维中，它变成相应单位的左乘映射，因而为同构。\cref{b2:lem:azumaya-local-toolkit}使 \(L_a\) 本身为同构。令 \(b=L_a^{-1}(1_A)\)，则 \(ab=1_A\)。为得到另一侧的逆，计算
\[
L_a(ba-1_A)=a(ba-1_A)=(ab)a-a=0.
\]
由 \(L_a\) 的单射性，\(ba=1_A\)，所以 \(a\) 可逆。
\end{proof}

% !TeX root = ../../Book2.tex
% 纲要标题：### C.3 环的 Brauer 群
% 纲要位置：工作记录/计划纲要.md，第 1820 行。

% 纲要要点（供目录核对）：
% * 环的 Brauer 群及其与域情形的区别
% #### C.3.1 等价关系与群运算
% #### C.3.2 分裂类、底环变换与域情形

\section{环的 Brauer 群}
\label{b2:sec:C03}

域的 Brauer 群把增加矩阵大小视为无关变化。一般交换环上的有限投射模可能没有统一的自由基，所以相应的可忽略因子应从矩阵代数扩展为 \(\operatorname{End}_R(P)\)，其中 \(P\) 忠实、有限生成且投射。上一节已经证明，这些因子确实属于 Azumaya 代数。

\subsection{等价关系与群运算}
\label{b2:subsec:C03:01}

\begin{proposition}\label{b2:prop:azumaya-tensor-opposite}
设 \(R\) 为非零交换含幺环。若 \(A,B\) 为 Azumaya \(R\) 代数，则 \(A^{\mathrm{op}}\) 与 \(A\otimes_RB\) 也为 Azumaya \(R\) 代数。
\end{proposition}
\begin{proof}
反代数不改变底模，张量积的忠实有限投射性由\cref{b2:lem:azumaya-projective-base-change}保证。对任意极大理想 \(\mathfrak m\)，这些代数的纤维分别为 \(A(\mathfrak m)^{\mathrm{op}}\) 与 \(A(\mathfrak m)\otimes_{\kappa(\mathfrak m)}B(\mathfrak m)\)。原纤维中心单；反代数保留单性与中心，张量积的中心单性由\cref{b2:thm:csa-tensor-product}保证。因此\cref{b2:thm:azumaya-fiber-criterion}适用于两个新代数。
\end{proof}

\begin{definition}\label{b2:def:ring-brauer-equivalence}
设 \(R\) 为非零交换含幺环、\(A,B\) 为 Azumaya \(R\) 代数。若存在忠实有限生成投射 \(R\) 模 \(P,Q\)，以及 \(R\) 代数同构
\begin{equation}\label{b2:eq:ring-brauer-equivalence}
A\otimes_R\operatorname{End}_R(P)
\cong B\otimes_R\operatorname{End}_R(Q),
\end{equation}
则称 \(A,B\) 在 \(R\) 上 Brauer 等价。
\end{definition}

此处同构必须保持指定的 \(R\) 标量，不能只要求抽象环同构。\chapref{b2:ch:18}已经定义域上的 Brauer 等价；这里使用同一名称，并将在下一个小节说明两种定义在域上重合。

\begin{theorem}\label{b2:thm:ring-brauer-group}
设 \(R\) 为非零交换含幺环。在 Azumaya \(R\) 代数上，\cref{b2:def:ring-brauer-equivalence}所定义的 Brauer 等价是等价关系。记 \([A]\) 为代数 \(A\) 的等价类。这些等价类在运算
\[
[A]+[B]=[A\otimes_RB]
\]
下构成交换群，零元为 \([R]\)，\([A]\) 的逆元为 \([A^{\mathrm{op}}]\)。这个群称为交换环 \(R\) 的 Brauer 群，记为 \(\operatorname{Br}(R)\)。
\end{theorem}
\symidx{\(\operatorname{Br}(R)\)：交换环的 Brauer 群}
\begin{proof}
取 \(P=Q=R\) 得自反性，交换同构两端得对称性。设
\[
A\otimes_R\operatorname{End}_R(P)\cong B\otimes_R\operatorname{End}_R(Q),
\qquad
B\otimes_R\operatorname{End}_R(U)\cong C\otimes_R\operatorname{End}_R(V).
\]
将第一式张量 \(\operatorname{End}_R(U)\)，第二式张量 \(\operatorname{End}_R(Q)\)，再交换中间因子，由式\eqref{b2:eq:azumaya-end-tensor}得到
\[
A\otimes_R\operatorname{End}_R(P\otimes_RU)
\cong C\otimes_R\operatorname{End}_R(V\otimes_RQ).
\]
\cref{b2:lem:azumaya-projective-base-change}保证这里的两个张量模仍忠实有限生成且投射，所以关系具有传递性。

若 \(A\) 与 \(A'\)、\(B\) 与 \(B'\) 分别通过式\eqref{b2:eq:ring-brauer-equivalence}等价，将两式张量并交换因子，再用同一个自同态张量公式，就得到 \(A\otimes_RB\) 与 \(A'\otimes_RB'\) 等价。因此运算不依赖代表。结合性、交换性和单位元分别由张量积的结合、交换和单位同构给出。

最后，\(A\) 自身是忠实有限生成投射模，定义中的同构给出
\[
A\otimes_RA^{\mathrm{op}}\cong\operatorname{End}_R(A).
\]
而 \(\operatorname{End}_R(A)\) 与 \(R\) 等价：在等价关系中分别选 \(R\) 和 \(A\) 作为附加模即可。所以 \([A]+[A^{\mathrm{op}}]=[R]\)，每个元素都有逆元。所有代数均以某个有限自由模的直和因子为底模，它们的同构类型构成集合，因而所述等价类确实构成一个群。
\end{proof}

特别地，\(\bigl[\operatorname{End}_R(P)\bigr]=0\)。这一步说明为何必须消去自同态代数，而不只消去固定大小的矩阵环。在不连通的基环上，\(P\) 的秩甚至可能在不同部分取不同的值。

\subsection{分裂类、底环变换与域情形}
\label{b2:subsec:C03:02}

设 \(R\) 为非零交换含幺环、\(A\) 为 Azumaya \(R\) 代数。若存在忠实有限生成投射 \(R\) 模 \(P\)，使 \(A\cong\operatorname{End}_R(P)\) 为 \(R\) 代数同构，则称 \(A\) 在 \(R\) 上分裂。\cref{b2:thm:ring-brauer-group}已经说明分裂代数代表零类。要证明反向，需从等价关系给出的附加自同态因子中恢复 \(A\)；下面的中心化子计算提供了这个恢复方法。

\begin{lemma}[自同态因子的中心化子]\label{b2:lem:azumaya-end-centralizer}
设 \(R\) 为非零交换含幺环、\(A\) 为结合含幺 \(R\) 代数，其结构同态的像位于中心。设 \(P\) 为忠实有限生成投射 \(R\) 模，令 \(E=\operatorname{End}_R(P)\)。在 \(A\otimes_RE\) 中记
\[
C=\{\,x\in A\otimes_RE\mid x(1\otimes e)=(1\otimes e)x\text{ 对全部 }e\in E\text{ 成立}\,\}.
\]
则 \(a\mapsto a\otimes\operatorname{id}_P\) 给出 \(R\) 代数同构 \(A\cong C\)。
\end{lemma}
\begin{proof}
上述映射的像包含在 \(C\) 中。由\cref{b2:lem:azumaya-projective-base-change}，\(E\) 有限生成。选其有限个 \(R\) 模生成元 \(e_1,\ldots,e_v\)，则 \(C\) 是映射
\[
A\otimes_RE\longrightarrow (A\otimes_RE)^v,
\qquad x\longmapsto\bigl(x(1\otimes e_i)-(1\otimes e_i)x\bigr)_{i=1}^v
\]
的核。局部化保持正合性，而且 \(e_i/1\) 生成每个局部化后的自同态代数，故取这个中心化子与局部化相容。

对任意极大理想 \(\mathfrak m\)，\cref{b2:lem:azumaya-local-toolkit,b2:lem:azumaya-projective-base-change}给出 \(P_{\mathfrak m}\cong R_{\mathfrak m}^n\)、\(n\ge1\)，以及
\[
(A\otimes_RE)_{\mathfrak m}\cong M_n(A_{\mathfrak m}).
\]
在这个识别下，与 \(1\otimes E\) 交换等价于与全部矩阵单位交换。与对角矩阵单位交换先使非对角条目为零，与其余矩阵单位交换再使对角条目相等。因此中心化子恰为 \(aI_n\)、\(a\in A_{\mathfrak m}\)，原映射的局部化就是 \(A_{\mathfrak m}\cong C_{\mathfrak m}\)。对原映射的核与余核应用\cref{b2:lem:azumaya-local-toolkit}的局部消失判据，得 \(A\cong C\)。映射公式保留乘法、单位和 \(R\) 标量，故它是 \(R\) 代数同构。
\end{proof}

\begin{proposition}\label{b2:prop:azumaya-zero-split}
设 \(R\) 为非零交换含幺环、\(A\) 为 Azumaya \(R\) 代数。\(A\) 代表 \(\operatorname{Br}(R)\) 中的零类，当且仅当存在忠实有限生成投射 \(R\) 模 \(T\)，使 \(A\cong\operatorname{End}_R(T)\) 为 \(R\) 代数同构。
\end{proposition}
\begin{proof}
若 \(A\cong\operatorname{End}_R(T)\)，\cref{b2:thm:ring-brauer-group}给出 \([A]=0\)。反过来，设 \([A]=0\)，按等价关系选忠实有限生成投射 \(R\) 模 \(P,Q\)，使
\[
A\otimes_RE\cong\operatorname{End}_R(Q),
\qquad E=\operatorname{End}_R(P).
\]
经这个同构，\(Q\) 成为左 \(E\) 模，且 \(E\) 中的 \(R\) 标量仍按原来的 \(R\) 模结构作用于 \(Q\)。

令 \(P^*=\operatorname{Hom}_R(P,R)\)，赋予 \(P\) 双模结构 \({}_EP_R\)，其作用为 \(ep=e(p)\)、\(pr=rp\)；赋予 \(P^*\) 双模结构 \({}_RP^*_E\)，其作用为 \((r\lambda)(p)=r\lambda(p)\)、\((\lambda e)(p)=\lambda(e(p))\)。对偶保留有限自由模的直和因子，所以 \(P^*\) 有限生成且投射，双对偶求值映射 \(P\rightarrow P^{**}\) 为同构。取式\eqref{b2:eq:azumaya-trace-unit}中的 \(p_i,\lambda_i\)；若 \(rP^*=0\)，则 \(r=\sum_i(r\lambda_i)(p_i)=0\)。因此 \(P^*\) 忠实，由\cref{b2:lem:azumaya-projective-base-change}还是生成元。取对偶给出 \(\operatorname{End}_R(P^*)\cong E^{\mathrm{op}}\)，将\cref{b2:cor:morita-inverse-bimodules}用于左 \(R\) 模 \(P^*\)，再用双对偶同构，得到双模同构
\[
{}_RP^*\otimes_EP_R\cong{}_RR_R,
\qquad
{}_EP\otimes_RP^*_E\cong{}_EE_E.
\]
因此 \(P\otimes_R-\) 与 \(P^*\otimes_E-\) 分别从左 \(R\) 模到左 \(E\) 模、从左 \(E\) 模到左 \(R\) 模，且由张量结合律互为拟逆。令 \(T=P^*\otimes_EQ\)，便有 \(Q\cong P\otimes_RT\)，以及由等价的全忠实性得到的 \(R\) 代数同构
\[
\operatorname{End}_E(Q)\cong\operatorname{End}_R(T).
\]
由\cref{b2:lem:azumaya-end-centralizer}，\(A\otimes_RE\) 中 \(1\otimes E\) 的中心化子为 \(A\)；在所给同构下，它恰好对应 \(\operatorname{End}_E(Q)\)。所以 \(A\cong\operatorname{End}_R(T)\)。

还须核对 \(T\) 的模性质。对 \(p\in P\)、\(\lambda\in P^*\)，记 \(\theta_{p,\lambda}\in E\) 为算子 \(x\mapsto\lambda(x)p\)。式\eqref{b2:eq:azumaya-trace-unit}给出右 \(E\) 模同态
\[
E^t\longrightarrow P^*,\qquad (e_i)\longmapsto\sum_i\lambda_i e_i,
\]
其截面为 \(\lambda\mapsto(\theta_{p_i,\lambda})_i\)，因为 \(\sum_i\lambda_i\theta_{p_i,\lambda}=\lambda\)。这些映射也保持 \(R\) 作用，取 \(-\otimes_EQ\) 后使 \(T\) 成为 \(Q^t\) 的 \(R\) 模直和因子。因此 \(T\) 有限生成且投射。若 \(rT=0\)，由 \(Q\cong P\otimes_RT\) 得 \(rQ=0\)，而 \(Q\) 忠实，故 \(r=0\)。这证明 \(T\) 忠实。
\end{proof}

零类因而恰好由分裂代数代表。不过，这些代数之间仍可能不同构；Brauer 等价消去了自同态因子，不能据此恢复原代数的大小。改变底环则会进一步改变哪些类能够分裂。

\begin{proposition}\label{b2:prop:ring-brauer-base-change}
设 \(R,S\) 为非零交换含幺环。环同态 \(f:R\rightarrow S\) 诱导群同态
\[
f^*:\operatorname{Br}(R)\longrightarrow\operatorname{Br}(S),
\qquad [A]\longmapsto[S\otimes_RA].
\]
这些同态保持恒等映射，并按底环变换次序保持复合。若 \(R=k\) 为域，则 \(\operatorname{Br}(k)\) 与\cref{b2:thm:brauer-group}构造的群相同。
\end{proposition}
\begin{proof}
\cref{b2:thm:azumaya-fiber-criterion}保证底环变换后的代数仍为 Azumaya 代数。对式\eqref{b2:eq:ring-brauer-equivalence}取 \(S\otimes_R-\)，式\eqref{b2:eq:azumaya-hom-base-change}把附加因子化为
\[
S\otimes_R\operatorname{End}_R(P)
\cong\operatorname{End}_S(S\otimes_RP),
\]
且 \(S\otimes_RP\) 仍忠实有限生成且投射。故等价关系保持。底环变换与张量积相容，并把 \(R\) 送到 \(S\)，所以所得映射为群同态；两次底环变换的结合识别给出复合相容性。

在域上，忠实有限生成投射模恰为非零有限维向量空间，其自同态代数恰为 \(M_n(k)\)、\(n\ge1\)。又由\cref{b2:prop:azumaya-field-matrix}，代数对象恰为中心单代数。因此等价关系变为 \(M_n(A)\cong M_m(B)\)，群运算也同为张量积，两个构造完全一致。
\end{proof}

例如，取 \(R=\mathbb R[t]\)，由 \(\mathbb R\rightarrow R\) 扩张得到 \(A=R\otimes_{\mathbb R}\mathbb H\)。这是 Azumaya \(R\) 代数，却不代表零类。若 \([A]=0\)，再按 \(t\mapsto0\) 底环变换到 \(\mathbb R\)，便有 \([\mathbb H]=0\)。\cref{b2:prop:brauer-division-representative}说明域上每个 Brauer 类有唯一的中心除代数代表，\(\mathbb H\ne\mathbb R\)，所以这不可能。另一方面，扩张到 \(\mathbb C[t]\) 后 \(A\) 变成 \(M_2\bigl(\mathbb C[t]\bigr)\)，其类随之变为零。

\ProseParagraphBreak
一般交换环上的 Brauer 群仍以代数及其模为对象，但域上“每类选一个除代数”的描述已不适用。例如 \(R=k\times k\) 的任何忠实代数都包含非零的中心幂等元 \((1,0)\)、\((0,1)\)，二者乘积为零，因而不可能是除环。环上的分类问题还须顾及各处的纤维如何在同一个有限投射模上相容。以自同态代数为零类、以张量积为加法的表述，也可参阅 Weibel 的《The K-book: An Introduction to Algebraic K-theory》\cite[Chapter II, Example 5.4.3]{Weibel2013KBook}；该处从对称幺半范畴的群化解释这组关系。

% !TeX root = ../../Book2.tex
% 纲要标题：### C.4 向量丛、自同态代数与 Morita 等价
% 纲要位置：工作记录/计划纲要.md，第 1828 行。

% 纲要要点（供目录核对）：
% * 有限投射模的局部坐标、自同态代数的粘合与 Morita 等价
% #### C.4.1 向量丛与自同态代数
% #### C.4.2 模范畴与几何信息
% ---

\section{向量丛、自同态代数与 Morita 等价}
\label{b2:sec:C04}

有限生成投射模在每个素理想处给出一个有限维纤维；要描述这些纤维如何组成同一个对象，还需主开覆盖上的自由模，以及交叠处满足相容条件的换基矩阵。这个局部坐标描述把有限生成投射模与向量丛联系起来，并使其自同态代数成为由共轭变换粘合的局部矩阵代数。Morita 等价则说明，从模范畴观察这些自同态代数时，哪些来自交换基环的信息得到保留。

\subsection{向量丛与自同态代数}
\label{b2:subsec:C04:01}

设 \(R\) 为非零交换含幺环、\(P\) 为有限生成投射 \(R\) 模。对 \(f\in R\)，主开集 \(D(f)=\{\mathfrak p\in\operatorname{Spec}R\mid f\notin\mathfrak p\}\) 上的模由 \(P_f=R[1/f]\otimes_RP\) 描述。\cref{b2:prop:projective-basic-open-basis}给出有限主开覆盖 \(D(f_i)\)，使每个 \(P_{f_i}\) 都自由。若局部秩均为 \(n\)，在这组覆盖中分别选基；交叠区域上的坐标由换基矩阵
\[
g_{ij}\in\operatorname{GL}_n(R_{f_if_j})
\]
变换，并在三重交叠上满足 \(g_{ij}g_{jk}=g_{ik}\)。这里约定 \([p]_i=g_{ij}[p]_j\)，两边均取到 \(R_{f_if_j}\) 上。这些主开集上的自由模及其相容换基矩阵，给出\term[xiangliang-cong]{向量丛}[vector bundle]的局部描述；本节只使用从已经给定的 \(P\) 导出的这一模型。不同连通部分可以有不同的局部秩。

\ProseParagraphBreak
对 \(E=\operatorname{End}_R(P)\)，同一组局部基给出 \(E_{f_i}\cong M_n(R_{f_i})\)。若 \(T\in E\) 的局部矩阵为 \(T_i\)，由坐标变换公式可得
\begin{equation}\label{b2:eq:endomorphism-bundle-transition}
T_i=g_{ij}T_jg_{ij}^{-1}.
\end{equation}
所以自同态代数使用共轭来粘合矩阵代数。这些公式保留加法、乘法与单位矩阵；因此局部矩阵的描述与全局代数结构相容。若 \(P\) 忠实，它给出\cref{b2:cor:azumaya-endomorphism}中的分裂 Azumaya 代数。

\ProseParagraphBreak
局部秩为一时，\(g_{ij}\) 是标量单位，局部自同态 \(T_j\) 也是标量，交换性使 \(g_{ij}T_jg_{ij}^{-1}=T_j\)。因此模本身的换基可能仍然非平凡，自同态代数的共轭换基却已经消失。是否存在全局模基，与自同态代数是否为矩阵环，是两个不同的问题。

\ProseParagraphBreak
\cref{b2:ex:krull-schmidt-nonunique}中的 \(R=\mathbb Z[s]\)、\(s^2=-5\)、\(I=(2,1+s)\) 给出同一局部秩下的例子。该处已证 \(I\) 非主且 \(I^2=2R\)。对 \(t\in I\)，映射 \(t\mapsto2t\) 及 \(t\mapsto(s-1)t/2\) 都取值于 \(R\)，后者使用 \(s-1\in I\) 和 \(I^2=2R\)。并有
\[
t=2(2t)+(1+s)\frac{(s-1)t}{2}.
\]
由\cref{b2:prop:finite-dual-basis}，\(I\) 是有限生成投射模。在 \(D(2)\) 上，\(I_2=R_2\)；在 \(D(3)\) 上，由 \(2=(1+s)(1-s)/3\)，有 \(I_3=(1+s)R_3\)。这两个主开集覆盖 \(\operatorname{Spec}R\)，所以局部秩处处为一；若有全局基，\(I\) 就应由一个元素生成，与非主性矛盾。

\ProseParagraphBreak
尽管 \(I\) 不自由，却有自然的代数同构 \(\operatorname{End}_R(I)\cong R\)。确实，取 \(0\ne t_0\in I\)，任一自同态 \(T\) 满足 \(t_0T(t)=tT(t_0)\)，所以它在分式域中是乘以某个 \(c\) 的映射。由 \(cI\subseteq I\)，有 \(cI^2\subseteq I^2=2R\)，消去 \(2\) 得 \(cR\subseteq R\)，故 \(c\in R\)。反之每个 \(c\in R\) 都给出这样的自同态，乘法和复合相容。这具体展示了秩一换基的共轭作用为何消失。

\ProseParagraphBreak
例如 \(R=k\times k\)，取 \(P=k\times k^2\)，由两个坐标分别作标量作用。这个模是 \(R^2\) 的直和因子且忠实，两个点上的秩分别为一、二，故不可能是某个固定秩的自由 \(R\) 模。其自同态代数为 \(k\times M_2(k)\)，是分裂 Azumaya \(R\) 代数，却不同构于任何 \(M_n(R)\)：在两个坐标上比较维数将同时要求 \(n^2=1\) 与 \(n^2=4\)。这正是环上的“分裂”比“一个固定大小的矩阵环”更自然的原因。

\subsection{模范畴与几何信息}
\label{b2:subsec:C04:02}

\begin{proposition}\label{b2:prop:split-azumaya-morita}
设 \(R\) 为非零交换含幺环、\(P\) 为忠实有限生成投射 \(R\) 模，令 \(E=\operatorname{End}_R(P)\)、\(P^*=\operatorname{Hom}_R(P,R)\)。赋予 \(P\) 双模结构 \({}_EP_R\)，赋予 \(P^*\) 双模结构 \({}_RP^*_E\)，其中
\[
ep=e(p),\qquad pr=rp,\qquad
(r\lambda)(p)=r\lambda(p),\qquad (\lambda e)(p)=\lambda(e(p)).
\]
则左模范畴之间的两张量函子
\[
\begin{aligned}
F:R\text{-}\mathbf{Mod}&\longrightarrow E\text{-}\mathbf{Mod},
&N&\longmapsto P\otimes_RN,\\
G:E\text{-}\mathbf{Mod}&\longrightarrow R\text{-}\mathbf{Mod},
&M&\longmapsto P^*\otimes_EM
\end{aligned}
\]
互为拟逆。\(G\) 也自然同构于 \(\operatorname{Hom}_E(P,-)\)，其中该 Hom 模的左 \(R\) 作用为 \((rf)(p)=f(pr)\)。
\end{proposition}
\begin{proof}
对有限投射模，对偶两次的求值映射 \(P\rightarrow P^{**}\) 是同构：这在有限自由模上由坐标直接成立，限制到直和因子后仍成立。因此取对偶给出
\[
\operatorname{End}_R(P^*)\cong E^{\mathrm{op}}.
\]
对偶模有限投射，而且在每个局部环上的秩与 \(P\) 相同，均为正。若 \(rP^*=0\)，在每个极大理想处用正秩自由模得 \(r/1=0\)，再由局部消失判据应用于 \(Rr\) 得 \(r=0\)。故 \(P^*\) 忠实，由\cref{b2:lem:azumaya-projective-base-change}还是生成元。

将\cref{b2:thm:morita-equivalence}用于左 \(R\) 模 \(P^*\)。该定理的自同态反环现在为 \(E\)，给出等价函子 \(\operatorname{Hom}_R(P^*,-)\) 及拟逆 \(P^*\otimes_E-\)。由\cref{b2:prop:finite-projective-hom-tensor}和双对偶同构，前一个函子自然同构于 \(P\otimes_R-\)，其中左 \(E\) 作用恰为求值作用。

等价的全忠实性又给出关于 \(N\) 自然的同构
\[
\operatorname{Hom}_E(P,P\otimes_RN)
\cong\operatorname{Hom}_R(R,N)\cong N.
\]
对任意左 \(E\) 模 \(M\)，已有自然同构 \(M\cong P\otimes_R(P^*\otimes_EM)\)。在上式取 \(N=P^*\otimes_EM\)，得到 \(\operatorname{Hom}_E(P,M)\cong P^*\otimes_EM\)，因此这个 Hom 函子也给出拟逆。
\end{proof}

以 \(R=k[t]\)、\(P=R^n\)、\(n\ge1\) 为例，取 \(a\in k\)、\(r\ge1\)。矩阵环的模范畴仍记录多项式变量 \(t\) 的作用。模 \(R/(t-a)\) 经上述等价成为 \(k^n\)，其中 \(t\) 作标量 \(a\)；模 \(R/(t-a)^r\) 则成为 \(\bigl(R/(t-a)^r\bigr)^n\)，保留 \((t-a)\) 的幂零层数。矩阵大小改变了底向量空间维数，却没有丢掉这些局部信息。

\ProseParagraphBreak
这与\cref{b2:prop:morita-center}相呼应：模范畴恒等函子的自然自变换环是原环的中心，因而 Morita 等价保持中心。对 Azumaya \(R\) 代数，中心由\cref{b2:prop:azumaya-center}确定为 \(R\)，所以模范畴仍保留这个交换基环。另一方面，中心相同不保证 Morita 等价；\(\mathbb R\) 与 \(\mathbb H\) 的中心相同，它们的唯一简单模的自同态除环分别为 \(\mathbb R\) 与 \(\mathbb H^{\mathrm{op}}\)，不可能由范畴等价相互对应。非零 Brauer 类正体现了这种中心之外的代数差别。

\ProseParagraphBreak
本节使用的几何语言都有具体的模论内容：主开集对应局部化，局部基对应向量丛的坐标，矩阵共轭对应自同态代数的粘合，模范畴等价对应 Morita 理论。一般非交换代数的几何还会研究分次模、商范畴及导出范畴等对象。进一步学习这些方向时，需要另行建立相应的几何或同调基础；不能仅凭一个非交换环的中心，就断言它具有某个完整的交换几何模型。


\begin{chaptersummary}
主开局部化使指定的底环元素可逆，让有限投射模获得局部自由基；剩余域纤维则进一步模去点处的极大理想，把局部模压为域上的向量空间。在底模忠实、有限生成且投射的前提下，左右乘包络映射是否同构，恰可由全部极大纤维是否中心单检验。一个泛纤维或稠密主开集上的矩阵模型不足以代替这个逐点条件。平展局部分裂还允许可分扩张的方向：\(\mathbb H\) 在 \(\mathbb R\) 的 Zariski 局部化中保持不变，却经 \(\mathbb R\rightarrow\mathbb C\) 分裂。

\ProseParagraphBreak
Azumaya 代数在任意非零交换底环变换下保持其定义条件，张量积与反代数也保持 Azumaya 性。环的 Brauer 群用忠实有限投射模的自同态代数作稳定因子：张量积给出加法，反代数给出逆元，Morita 理论说明零类恰由分裂代数 \(\operatorname{End}_R(P)\) 代表。域上这恢复矩阵稳定等价，环上则必须容纳未必有全局基、局部秩可能变化的 \(P\)。零类在每个剩余域纤维上仍为零，因而一个非分裂纤维已能排除全局零类。

\ProseParagraphBreak
向量坐标按可逆矩阵变换，自同态代数的坐标按矩阵共轭变换，标量换基在共轭中消失。秩一投射模 \(L\) 即使不自由，仍有 \(\operatorname{End}_R(L)\cong R\)；一般的 \(P\) 换成 \(P\otimes_RL\) 也不改变其自同态代数。分裂 Azumaya 代数与基环有等价的模范畴，这种等价保留中心及模中的幂零厚度，而不保留矩阵大小。中心相同的 \(\mathbb R\) 与 \(\mathbb H\) 却不 Morita 等价，其简单对象的自同态除环仍能区分二者。
\end{chaptersummary}

\BookExercises

\begin{exercise}\label{b2:ex:C-transpose-tensor}
设 \(A=M_n(R)\)。在 \(A\otimes_RA^{\mathrm{op}}\) 中分别找出转置算子 \(X\mapsto X^{\mathsf T}\) 与算子 \(X\mapsto\operatorname{Tr}(X)I_n\) 在 \(\mu_A\) 下的原像。说明当 \(n\ge2\) 时，转置不保持乘法仍不妨碍它属于 \(\mu_A\) 的像。
\end{exercise}
\begin{solution}
由\cref{b2:prop:azumaya-field-matrix}证明中的矩阵单位公式，两个原像分别为
\[
\sum_{i,j}E_{ij}\otimes E_{ij}^{\mathrm{op}},
\qquad
\sum_{i,j}E_{ij}\otimes E_{ji}^{\mathrm{op}}.
\]
第一个和把 \(E_{pq}\) 送到 \(E_{qp}\)；第二个和在 \(p\ne q\) 时给零，在 \(p=q\) 时给 \(\sum_iE_{ii}=I_n\)。这就分别给出所求算子。\(\mu_A\) 的目标是全部 \(R\) 线性自同态，而不是代数自同态的集合；乘法在目标中表示算子复合，所以不要求每个像算子自身保持矩阵乘法。
\end{solution}

\begin{exercise}\label{b2:ex:C-dual-number-failure}
设 \(k\) 为域，\(A=k[\varepsilon]/(\varepsilon^2)\)。直接找出 \(\mu_A\) 的非零核元素，并说明忠实有限自由的底模不能单独保证 Azumaya 性。
\end{exercise}
\begin{solution}
因为 \(A\) 交换，\(z=\varepsilon\otimes1-1\otimes\varepsilon\) 作用于任意 \(x\) 都给出 \(\varepsilon x-x\varepsilon=0\)。张量基 \(1\otimes1,\varepsilon\otimes1,1\otimes\varepsilon,\varepsilon\otimes\varepsilon\) 说明 \(z\ne0\)，即使特征为二也如此。\(A\) 的底向量空间维数为二，因而忠实且有限自由，但 \(\mu_A\) 不单射，所以不是 Azumaya 代数。
\end{solution}

\begin{exercise}\label{b2:ex:C-generic-fiber-insufficient}
设 \(k\) 为域、\(R=k[t]\)，并取矩阵子代数
\[
B=\left\{\begin{pmatrix}a&tb\\tc&d\end{pmatrix}:a,b,c,d\in R\right\}.
\]
比较 \(B[1/t]\)、纤维 \(B/tB\) 及外在矩阵条目求值 \(t=0\) 的像。求后两者之间自然映射的核，并说明为何该映射不能用来检验 Azumaya 性。
\end{exercise}
\begin{solution}
\(e=E_{11},f=E_{22},x=tE_{12},y=tE_{21}\) 是 \(B\) 的 \(R\) 基，所以底模有限自由且忠实。逆置 \(t\) 后，这组基产生全部矩阵单位，故 \(B[1/t]=M_2(k[t,t^{-1}])\)。在 \(B/tB\) 中四个基元素仍有线性无关的像，关系为
\[
\bar x^2=\bar y^2=\bar x\bar y=\bar y\bar x=0,
\qquad \bar e\bar x=\bar x=\bar x\bar f,
\quad \bar f\bar y=\bar y=\bar y\bar e.
\]
因此 \(J=k\bar x+k\bar y\) 是非零平方零双边理想，纤维不是中心单代数。纤维判据排除 \(B\) 的 Azumaya 性。

逐条目令 \(t=0\) 却只得到对角代数 \(k\times k\)，并将 \(\bar x,\bar y\) 都送到零。这个求值通过 \(B/tB\) 给出满同态，其核恰为 \(J\)，因为它在 \(k\bar e\oplus k\bar f\) 上单射。这里外在求值丢掉了两个仍属于内在纤维的基方向；\cref{b2:prop:C-generic-fiber-insufficient}中的单侧倍数与此处双侧倍数，都须先按底模取商，而不是把矩阵嵌入在求值后仍看作单射。
\end{solution}

\begin{exercise}\label{b2:ex:C-finite-fiber-hypothesis}
计算 \(M=\mathbb Z[1/2]\) 作为 \(\mathbb Z\) 模的全部极大纤维，以及它在极大理想 \((2)\) 处的局部化。判断它是否忠实、平坦及有限生成。再设 \(k\) 为域，取 \(R=k[\varepsilon]/(\varepsilon^3)\)，计算商映射 \(u:R\rightarrow R/(\varepsilon^2)\) 的纤维与核，判断目标模是否投射。
\end{exercise}
\begin{solution}
对奇素数 \(p\)，二在 \(\mathbb F_p\) 中可逆，故 \(M/pM\cong\mathbb F_p\)；对 \(p=2\)，有 \(2M=M\)，所以纤维为零。但 \(M_{(2)}=\mathbb Q\)：局部化先后逆置了二及全部奇整数。于是同一点的局部化非零，纤维却为零。

\(M\subseteq\mathbb Q\) 的整数作用忠实，局部化又由\cref{b2:prop:localization-flat}保证平坦。若它由有限个分数生成，可取 \(2^N\) 同时清去这些分母，它们的整数线性组合就都属于 \(2^{-N}\mathbb Z\)，不能生成 \(2^{-N-1}\)。所以它不有限生成。

第二个环的唯一极大理想为 \((\varepsilon)\)，商映射在剩余域上为 \(k\) 的恒等映射，但原核为非零的 \(k\varepsilon^2\)。若目标投射，商映射有截面；一的像只能为 \(1+c\varepsilon^2\)，且须被 \(\varepsilon^2\) 零化，而 \(\varepsilon^2(1+c\varepsilon^2)=\varepsilon^2\ne0\)，矛盾。这个较高幂零次数的例子仍显示：纤维同构不使任意满射分裂。
\end{solution}

\begin{exercise}\label{b2:ex:C-faithful-projective-rank}
设 \(E=(e_{ij})\in M_N(R)\) 满足 \(E^2=E\)，令 \(P=E R^N\)。证明所有 \(\lambda(p)\)、\(\lambda\in P^*,p\in P\) 生成的理想，等于矩阵条目 \(e_{ij}\) 生成的理想；由此判断 \(P\) 忠实的充要条件。对 \(R=k\times k\)、\(E=\operatorname{diag}\bigl((1,1),(0,1)\bigr)\) 作具体计算。
\end{exercise}
\begin{solution}
\(R^N=P\oplus\ker E\)，所以 \(P\) 有限生成且投射。第 \(i\) 个坐标泛函在 \(Ee_j\) 处取值为 \(e_{ij}\)，故条目生成的理想包含于这些取值生成的理想。反过来，\(\lambda\) 可沿投影 \(E\) 延伸到 \(R^N\)，写成某个行向量 \(c\)；对 \(p=Ea\)，有 \(\lambda(p)=cEa\)，属于条目生成的理想，所以两理想相等。

若 \(P\) 忠实，\cref{b2:lem:azumaya-projective-base-change}中的单位元等式使这个理想为 \(R\)。若该理想为 \(R\)，而 \(rP=0\)，则 \(r e_{ij}=0\) 对全部条目成立，故 \(rR=0\)，即 \(r=0\)。因此忠实性恰等价于条目生成单位理想。在给定矩阵中条目包含 \((1,1)=1_R\)，故 \(P\) 忠实；两个因子上的投影矩阵分别为 \(\operatorname{diag}(1,0)\) 与 \(I_2\)，局部秩分别为一、二。
\end{solution}

\begin{exercise}\label{b2:ex:C-dual-endomorphism}
设 \(P\) 为有限生成投射 \(R\) 模。写出 \(\operatorname{End}_R(P)^{\mathrm{op}}\cong\operatorname{End}_R(P^*)\) 的映射，并说明有限投射性在满射性中的作用。
\end{exercise}
\begin{solution}
将 \(T\) 送到 \(T^*: \lambda\mapsto\lambda\circ T\)。对复合有 \((TU)^*=U^*T^*\)，所以这是反环到自同态环的同态。有限投射性保证求值映射 \(j:P\rightarrow P^{**}\) 是同构。对 \(V\in\operatorname{End}_R(P^*)\)，定义 \(T=j^{-1}V^*j\)。在 \(p,\lambda\) 处计算得
\[
\lambda\bigl(T(p)\bigr)=V(\lambda)(p),
\]
故 \(T^*=V\)。同一等式也给出单射性。没有双对偶同构时，这种把对偶上的任意算子拉回 \(P\) 的步骤未必成立。
\end{solution}

\begin{exercise}\label{b2:ex:C-integral-quaternions}
在 \(\mathbb H\) 中取整数系数四元数环 \(B=\mathbb Z1+\mathbb Zi+\mathbb Zj+\mathbb Zij\)。证明 \(B\) 不是 \(\mathbb Z\) 上的 Azumaya 代数，但 \(B[1/2]\) 是 \(\mathbb Z[1/2]\) 上的 Azumaya 代数。
\end{exercise}
\begin{solution}
\(B\) 为秩四的忠实自由模。在模二纤维中，\(i,j\) 交换，令 \(u=\bar i+1,v=\bar j+1\)，则 \(u^2=v^2=0\)。四个基元素给出同构
\[
B/2B\cong\mathbb F_2[u,v]/(u^2,v^2).
\]
这个四维交换代数不是以 \(\mathbb F_2\) 为中心的单代数，所以纤维判据排除第一项。\(\mathbb Z[1/2]\) 的极大理想则恰来自奇素数 \(p\)，相应纤维为 \((-1,-1)_{\mathbb F_p}\)。\cref{b2:prop:quaternion-basis-splitting}说明在特征非二、两个参数非零时，这个四元数代数中心单。故全部极大纤维中心单，\cref{b2:thm:azumaya-fiber-criterion}给出第二项。这里使二可逆，正是从基谱中去掉破坏中心单性的模二纤维；其余各点的纤维已经满足判据。
\end{solution}

\begin{exercise}\label{b2:ex:C-unit-fibers}
令 \(R=\mathbb Z/4\mathbb Z\)、\(A=M_2(R)\)。对每个 \(b\in R\)，判断
\[
X_b=\begin{pmatrix}1&b\\b&1\end{pmatrix},\qquad
Y=\begin{pmatrix}2&1\\1&2\end{pmatrix}
\]
是否可逆，并求可逆时的逆。对不可逆的 \(X_b\) 给出一个非零核向量，核对纤维判据与直接计算是否一致。
\end{exercise}
\begin{solution}
\(R\) 只有极大理想 \((2)\)，其剩余域为 \(\mathbb F_2\)。模二以后，\(X_b\) 在 \(b\) 为偶数时成为单位矩阵，在 \(b\) 为奇数时成为全一矩阵。由\cref{b2:prop:C-unit-fiber-criterion}，\(X_b\) 恰在 \(b=0,2\) 时可逆。这时 \(b^2=0\) 于 \(R\)，直接相乘得到
\[
X_b^{-1}=\begin{pmatrix}1&-b\\-b&1\end{pmatrix}.
\]
若 \(b=1,3\)，则 \(b^2=1\)，非零列向量 \((-b,1)^{\mathsf T}\) 属于 \(X_b\) 的核，确实不能可逆。

\(Y\) 的模二纤维是交换两个坐标的可逆矩阵，故原矩阵也可逆。又 \(Y^2=\begin{psmallmatrix}5&4\\4&5\end{psmallmatrix}=I_2\) 于 \(R\)，所以 \(Y^{-1}=Y\)。这里先在唯一剩余域中判定，再在原零因子环中验证双侧逆，并未把“行列式非零”当作环上可逆的条件。
\end{solution}

\begin{exercise}\label{b2:ex:C-product-brauer}
设 \(R=R_1\times R_2\)，两个因子均为非零交换环。证明
\[
\operatorname{Br}(R)\cong\operatorname{Br}(R_1)\times\operatorname{Br}(R_2).
\]
\end{exercise}
\begin{solution}
令 \(e_1=(1,0),e_2=(0,1)\)。每个 \(R\) 模分解为 \(e_1M\oplus e_2M\)，每个 \(R\) 代数分解为 \(e_1A\times e_2A\)。有限投射性由直和因子刻画逐因子等价；忠实性也逐因子等价。不同因子之间的张量积为零，因为 \(e_1e_2=0\)。同态必须保留这两个中心幂等元的作用，所以自同态代数亦逐因子分解，\(\mu_A\) 因而分解为两个包络映射。由此 Azumaya 性逐因子等价。

Brauer 等价中的模 \(P,Q\) 同样逐因子分解，故给出两边的等价蕴含各分量等价。反向，把两个分量上的附加模分别取积，就得到 \(R\) 上的忠实有限投射模及原等价所需同构。最后张量积逐因子计算，所以这个等价类的双射保持群运算。
\end{solution}

\begin{exercise}\label{b2:ex:C-polynomial-brauer-injection}
证明常数嵌入 \(R\rightarrow R[t]\) 诱导的 \(\operatorname{Br}(R)\rightarrow\operatorname{Br}\bigl(R[t]\bigr)\) 为单射。这个论证是否已经证明它为满射？
\end{exercise}
\begin{solution}
取求值同态 \(R[t]\rightarrow R\)、\(t\mapsto0\)，它与常数嵌入的复合为 \(R\) 上的恒等映射。由\cref{b2:prop:ring-brauer-base-change}的复合相容性，对应 Brauer 群同态的复合也是恒等映射。因此常数嵌入诱导的同态有左逆，必单射。存在左逆不蕴含满射；这里没有说明任意 \(R[t]\) 上的 Azumaya 代数类是否由常数底环变换得到。
\end{solution}

\begin{exercise}\label{b2:ex:C-quaternion-laurent}
令 \(R=\mathbb R[t,t^{-1}]\)、\(A=R\otimes_{\mathbb R}\mathbb H\)。证明 \([A]\ne0\)，但 \(2[A]=0\)，并在扩张到 \(\mathbb C[t,t^{-1}]\) 后写出一个矩阵模型。
\end{exercise}
\begin{solution}
按 \(t\mapsto1\) 底环变换得到 \(\mathbb H\)，其除代数代表非平凡，所以 \([A]\) 不为零。四元数共轭式\eqref{b2:eq:quaternion-conjugation}给出 \(A\cong A^{\mathrm{op}}\)，故\cref{b2:thm:ring-brauer-group}使 \(2[A]=0\)。复系数下取
\[
i\longmapsto\begin{pmatrix}\mathrm i&0\\0&-\mathrm i\end{pmatrix},
\qquad
j\longmapsto\begin{pmatrix}0&-1\\1&0\end{pmatrix}.
\]
这两个矩阵平方均为 \(-I\) 且反交换。\(I\) 与它们及其乘积线性无关，并构成四维自由模的基；因此给出与 \(M_2\bigl(\mathbb C[t,t^{-1}]\bigr)\) 的代数同构。
\end{solution}

\begin{exercise}\label{b2:ex:C-line-twist}
沿用本附录的 \(R=\mathbb Z[s]\)、\(s^2=-5\)、\(I=(2,1+s)\)。证明乘法映射 \(I\otimes_RI\rightarrow2R\) 为同构，并写出生成元 \(2\) 的一个张量原像。对 \(P=R\oplus I\)，构造 \(P\otimes_RI\cong P\)，说明非自由秩一投射模张量变换并不一定改变一个较高秩模的同构类型。
\end{exercise}
\begin{solution}
已知 \(I^2=2R\)，故乘法映射满射。\(I\otimes_RI\) 是有限生成投射模，局部秩为一；\(2R\cong R\)、\(r\mapsto2r\)，因为 \(R\) 是整环。每个极大纤维上，乘法因满射成为一维空间之间的同构，由\cref{b2:lem:azumaya-local-toolkit}，原映射为同构。一个明确原像为
\[
z=2\otimes(1+s)-(1+s)\otimes(1+s)-2\otimes2,
\]
因为其乘积为 \(2(1+s)-(1+s)^2-4=2\)。于是 \(2r\mapsto rz\) 给出逆映射。

张量对直和分配，所得同构可以写为
\[
P\otimes_RI\cong I\oplus(I\otimes_RI)
\cong I\oplus2R\cong I\oplus R\cong R\oplus I=P.
\]
其中第二个映射是上述乘法，第三个把 \(2r\) 送到 \(r\)，最后交换两个直和因子。\(I\) 仍因非主而不自由，但它对这个 \(P\) 的张量变换已可由模同构消去。这比\cref{b2:prop:C-line-twist-invariance}所保证的自同态代数不变更强，是当前具体理想的额外性质。
\end{solution}

\begin{exercise}\label{b2:ex:C-base-ring-isomorphism}
令 \(k\) 为域、\(R=k\times k\times k\)，并按坐标作用取 \(P=k\times k^2\times k^3\)、\(Q=k^2\times k^3\times k\)。求两个自同态代数，证明它们抽象同构且 Brauer 等价，却非 \(R\) 代数同构。若 \(\tau(r_1,r_2,r_3)=(r_2,r_3,r_1)\)，写出一个满足 \(\Phi(ra)=\tau(r)\Phi(a)\) 的环同构，并说明改变目标的指定标量作用后如何使它成为 \(R\) 代数同构。
\end{exercise}
\begin{solution}
这两个模均为某个有限自由 \(R\) 模的直和因子，且三个分量都非零，故忠实有限投射。其自同态代数分别为
\[
A=k\times M_2(k)\times M_3(k),\qquad
B=M_2(k)\times M_3(k)\times k,
\]
所以都代表零 Brauer 类。映射 \(\Phi(a_1,a_2,a_3)=(a_2,a_3,a_1)\) 为抽象环同构，并逐分量满足 \(\Phi(ra)=\tau(r)\Phi(a)\)。它不是保留原标量的同构；任何 \(R\) 代数同构都须固定第一中央幂等元的作用，因而要给出 \(k\cong M_2(k)\) 的 \(k\) 代数同构，维数一与四使之不可能。

若把 \(B\) 的结构同态改为 \(r\mapsto\tau(r)1_B\)，记所得 \(R\) 代数为 \(B_\tau\)，上述等式就说明 \(\Phi:A\rightarrow B_\tau\) 是 \(R\) 代数同构。同一个抽象环配上不同的标量嵌入，产生了不同的 \(R\) 代数问题；这也给出\cref{b2:prop:C-base-algebra-structure}之外的显式半线性坐标变换。
\end{solution}

\begin{exercise}\label{b2:ex:C-transition-conjugation}
在 \(R=k[z,z^{-1}]\) 上取两个自由模坐标，其变换矩阵为 \(g=\operatorname{diag}(z,1)\)。若一个算子的第二组坐标矩阵为 \(T=\begin{pmatrix}a&b\\c&d\end{pmatrix}\)，求第一组坐标矩阵。再说明把坐标变换乘以任意可逆标量，并不改变自同态代数的变换规则。
\end{exercise}
\begin{solution}
由式\eqref{b2:eq:endomorphism-bundle-transition}，第一组坐标矩阵为
\[
gTg^{-1}=\begin{pmatrix}a&zb\\z^{-1}c&d\end{pmatrix}.
\]
若 \(u\in R^\times\)，则 \((ug)T(ug)^{-1}=u(gTg^{-1})u^{-1}=gTg^{-1}\)，因为标量矩阵与全部矩阵交换。这是向量坐标的变化与自同态坐标的变化所保留信息不同的一个具体表现。
\end{solution}

\begin{exercise}\label{b2:ex:C-matrix-fat-point}
令 \(R=k[t]\)、\(E=M_n(R)\)、\(r\ge1\)。在矩阵 Morita 等价下求 \(R/(t^r)\) 的像，计算其 \(k\) 维数、合成长度及自同态环，并说明 \(t\) 的幂零指数。
\end{exercise}
\begin{solution}
像为列模 \(M=\bigl(R/(t^r)\bigr)^n\)，其 \(k\) 维数为 \(nr\)。\(R/(t^r)\) 的链
\[
0\subset(t^{r-1})/(t^r)\subset\cdots\subset(t)/(t^r)\subset R/(t^r)
\]
有 \(r\) 个相邻简单商，均为 \(k\)。\cref{b2:prop:morita-invariants}保持合成长度，故 \(M\) 的长度也是 \(r\)。全忠实性给出
\[
\operatorname{End}_E(M)\cong\operatorname{End}_R\bigl(R/(t^r)\bigr)\cong R/(t^r),
\]
最后一个同构在一处求值，任意像均由乘该元素实现。\(t^rM=0\)，而 \(t^{r-1}\) 不零化第一坐标的单位向量，所以幂零指数恰为 \(r\)。
\end{solution}

\begin{exercise}\label{b2:ex:C-commutative-azumaya}
证明一个交换 \(R\) 代数 \(A\) 若为 Azumaya 代数，其结构同态 \(R\rightarrow A\) 必为同构。再解释一个次数大于一的有限可分域扩张为何不构成原底域上的 Azumaya 代数。
\end{exercise}
\begin{solution}
\(A\) 交换使 \(Z(A)=A\)，而\cref{b2:prop:azumaya-center}给出 \(Z(A)=R1_A\)。结构同态因忠实性单射，又由该等式满射，故为同构。若 \(L/k\) 为所述扩张，\(L\) 虽然是域并且在 \(k\) 上有限维，但中心为 \(L\ne k\)，不满足中心单条件。可分性不代替中心必须等于底域这一要求。
\end{solution}

\begin{exercise}\label{b2:ex:C-fiber-brauer-obstruction}
取 \(R=\mathbb R\times\mathbb R\)、\(A=\mathbb H\times\mathbb H\)，以及逐分量包含给出的
\[
R\longrightarrow S=\mathbb C\times\mathbb R
\longrightarrow T=\mathbb C\times\mathbb C.
\]
核对 \(A\) 为 \(R\) 上底模秩四的 Azumaya 代数。计算其 Brauer 类的阶，以及在 \(S,T\) 上的像的阶，说明只使一个分量的纤维分裂为何不足以使整个类为零。
\end{exercise}
\begin{solution}
两分量均为四维实中心单代数；分别选基便得 \(A\cong R^4\) 作为 \(R\) 模。两个极大纤维都是 \(\mathbb H\)，所以纤维判据给出 Azumaya 性。投影到任一分量把 \([A]\) 送到非零的 \([\mathbb H]\)，由\cref{b2:prop:ring-brauer-base-change}，\([A]\ne0\)。逐分量的四元数共轭给出 \(A\cong A^{\mathrm{op}}\)，故 \(2[A]=0\)，其阶恰为二。

第一次底环变换给出 \(S\otimes_RA\cong M_2(\mathbb C)\times\mathbb H\)。第二个分量仍有非零 Brauer 类，所以整个 \(S\) 上的类非零；同一反代数同构使其阶仍为二。第二次底环变换则给出 \(T\otimes_RA\cong M_2(\mathbb C)\times M_2(\mathbb C)=M_2(T)\)，于是类为零，阶为一。一个分量上的障碍虽已消失，另一个分量仍能阻止全局分裂。
\end{solution}

\begin{exercise}\label{b2:ex:C-morita-center-not-enough}
设 \(k\) 为域，\(T\) 为二阶上三角 \(k\) 矩阵环。证明 \(Z(T)=kI_2\)，但 \(T\) 既不是 \(k\) 上的 Azumaya 代数，也不与 \(k\) 在 Morita 意义下等价。
\end{exercise}
\begin{solution}
与 \(E_{11},E_{22}\) 交换使中心矩阵为对角矩阵，再与 \(E_{12}\) 交换使两个对角元相同，故中心恰为标量。\(kE_{12}\) 是非零真双边理想，所以 \(T\) 不中心单，由\cref{b2:prop:azumaya-field-matrix}也就不是 Azumaya 代数。

取一维左模 \(S_1,S_2\)，分别让三角矩阵按左上角、右下角条目作用。它们均简单且不同构，因为 \(E_{11}\) 在前者上为恒等、在后者上为零。而 \(k\) 只有一种简单模的同构类型。\cref{b2:prop:morita-invariants}说明等价保留简单模类型，故两个模范畴不等价。中心同构并未消除三角代数中的理想及简单模差别。
\end{solution}
